\documentclass[11pt,a4paper]{article}

\usepackage[margin=2cm]{geometry}
\usepackage{amssymb}
\usepackage{amsmath}
\usepackage{siunitx}
\usepackage{multicol}
\usepackage{enumitem}
\setlength{\parskip}{0.6em}
\setlength{\parindent}{0pt}
\sisetup{per-mode=symbol}



\begin{document}

\begin{center}
    {\LARGE \textbf{Energy Resources Revision Booklet}}\\[0.2cm]
    {\large GCSE Physics -- Year 9}
\end{center}

\hrule
\vspace{0.5cm}

\textbf{Instructions}

\begin{itemize}[noitemsep]
    \item Work through all the tasks in order.
    \item Show full working for all calculations.
    \item Use your Collins AQA Physics Student Book (pages 32--35) \emph{if you want}, but everything you need is in this booklet.
    \item Aim to spend about 2--3 hours on this booklet.
\end{itemize}

\hrule
\vspace{0.5cm}

\section*{1. Quick Recap: Energy, Power and Efficiency}

\textbf{Key ideas}

\begin{itemize}[noitemsep]
    \item \textbf{Energy} is measured in joules (J).
    \item \textbf{Power} is the rate of energy transfer:
    \[
        P = \frac{E}{t}
    \]
    where $P$ is power in watts (W), $E$ is energy in joules (J), and $t$ is time in seconds (s).
    \item We can also write:
    \[
        E = P \times t
    \]
    \item \textbf{Efficiency} tells you how much of the input energy is usefully transferred:
    \[
        \text{Efficiency} = \frac{\text{useful output energy}}{\text{total input energy}}
    \]
    or
    \[
        \text{Efficiency (\%)} = \frac{\text{useful output}}{\text{total input}} \times 100\%
    \]
\end{itemize}

\subsection*{Task 1: Key terms}

Fill in the blanks using the words in the box.

\textbf{Word box:} efficiency, power, joules, wasted, useful, seconds, watts

\begin{enumerate}[label=\alph*)]
    \item Energy is measured in \underline{\hspace{3cm}}.
    \item Time is measured in \underline{\hspace{3cm}}.
    \item Power is measured in \underline{\hspace{3cm}}.
    \item Power is the rate at which energy is transferred per unit time: it tells us how quickly energy is used or transferred. A high power device transfers energy \underline{\hspace{2cm}}.
    \item \underline{\hspace{3cm}} energy is the energy that does the job we want.
    \item The rest of the energy is \underline{\hspace{3cm}} energy (usually heat or sound).
    \item \underline{\hspace{3cm}} tells us how good a device is at turning input energy into useful output energy.
\end{enumerate}

\subsection*{Task 2: Power calculations}

Use $E = P \times t$. Remember to convert minutes to seconds.

\begin{enumerate}[label=\arabic*)]
    \item A phone charger has a power of \SI{10}{\watt}. It charges a phone for \SI{30}{\minute}. \\
    Calculate the energy transferred to the phone in joules.
    \vspace{2cm}

    \item A kettle has a power of \SI{2.0}{\kilo\watt}. It is on for \SI{90}{\second}. \\
    (a) Convert \SI{2.0}{\kilo\watt} to watts. \\
    (b) Calculate the energy transferred.
    \vspace{3cm}

    \item A hairdryer transfers \SI{72 000}{\joule} of energy in \SI{60}{\second}. \\
    Calculate its power.
    \vspace{2.5cm}

    \item A low-energy light bulb uses \SI{12}{\watt} of power. It is left on for \SI{4}{\hour}. \\
    (a) Convert \SI{4}{\hour} into seconds. \\
    (b) Calculate the energy it uses in joules.
    \vspace{3cm}
\end{enumerate}
\newpage
\subsection*{Task 3: Efficiency calculations}

\begin{enumerate}[label=\arabic*)]
    \item A bulb is supplied with \SI{100}{\joule} of electrical energy. It gives out \SI{20}{\joule} of light energy, the rest is heat. \\
    (a) Calculate the efficiency as a decimal. \\
    (b) Calculate the efficiency as a percentage.
    \vspace{3cm}

    \item A motor is \SI{40}{\percent} efficient and has an input power of \SI{200}{\watt}. \\
    Calculate the useful output power.
    \vspace{2.5cm}

    \item A device has an input energy of \SI{500}{\joule} and an efficiency of 0.6. \\
    (a) Calculate the useful output energy. \\
    (b) Calculate the wasted energy.
    \vspace{3cm}

    \item \textbf{Challenge:} A heater transfers \SI{18 000}{\joule} of useful energy in \SI{2}{\minute}. Its efficiency is \SI{90}{\percent}. \\
    Calculate the \emph{total} input energy.
    \vspace{3cm}
\end{enumerate}

\newpage

\section*{2. Energy Resources Overview}

\textbf{Key ideas}

\begin{itemize}[noitemsep]
    \item An \textbf{energy resource} is something we can use to generate electricity or provide energy for heating and transport.
    \item \textbf{Renewable} resources can be replaced naturally and will not run out on human timescales (e.g.\ wind, solar, hydro).
    \item \textbf{Non-renewable} resources will run out; they are used faster than they are formed (e.g.\ fossil fuels, nuclear fuel).
    \item Important issues: reliability, start-up time, cost, environmental impact, and whether they can meet base load and peak demand.
\end{itemize}

\subsection*{Task 4: Sorting energy resources}

\textbf{(a)} Put each resource in the correct list.

\textbf{Resource list:} coal, oil, natural gas, nuclear, wind, solar, hydroelectric, tidal, wave, geothermal, biomass

\begin{multicols}{2}
\textbf{Renewable:}
\begin{itemize}
    \item 
    \item 
    \item 
    \item 
    \item 
    \item 
\end{itemize}

\columnbreak

\textbf{Non-renewable:}
\begin{itemize}
    \item 
    \item 
    \item 
    \item 
\end{itemize}
\end{multicols}

\vspace{0.5cm}

\textbf{(b)} Explain in your own words what \textbf{renewable} means.
\vspace{2cm}

\textbf{(c)} Explain in your own words what \textbf{non-renewable} means.
\vspace{2cm}
\newpage
\subsection*{Task 5: How do they work?}

For each of the following energy resources, write 1--2 sentences to describe \emph{how} it is used to generate electricity.

\begin{enumerate}[label=\arabic*)]
    \item Coal-fired power station
    \vspace{1.7cm}
    \item Nuclear power station
    \vspace{1.7cm}
    \item Wind turbine
    \vspace{1.7cm}
    \item Solar panels (solar cells)
    \vspace{1.7cm}
    \item Hydroelectric power station
    \vspace{1.7cm}
    \item Tidal power
    \vspace{1.7cm}
\end{enumerate}

\newpage

\subsection*{Task 6: Comparing energy resources}

Complete as much of the table as you can. Use ticks ($\checkmark$) and crosses ($\times$) or short notes.

\begin{center}
\begin{tabular}{|l|c|c|c|c|c|}
\hline
\textbf{Resource} & \textbf{Renewable?} & \textbf{Reliable?} & \textbf{Slow/fast} & \textbf{CO$_2$?} & \textbf{Key issues} \\
\hline
Coal & & & & & \\
\hline
Natural gas & & & & & \\
\hline
Nuclear & & & & & \\
\hline
Wind & & & & & \\
\hline
Solar & & & & & \\
\hline
Hydroelectric & & & & & \\
\hline
Tidal & & & & & \\
\hline
Biomass & & & & & \\
\hline
\end{tabular}
\end{center}

Hints:
\begin{itemize}[noitemsep]
    \item ``Reliable'' means it can supply power whenever we need it.
    \item ``Slow/fast'' means the start-up time (how quickly it can be switched on).
    \item CO$_2$ means whether it releases carbon dioxide during normal operation.
\end{itemize}

\vspace{0.5cm}

\subsection*{Task 7: Advantages and disadvantages}

For each resource below, give \textbf{one advantage} and \textbf{one disadvantage}.

\begin{enumerate}[label=\arabic*)]
    \item Nuclear power
    \begin{itemize}
        \item Advantage:
        \vspace{5cm}
        \item Disadvantage:
        \vspace{6cm}
    \end{itemize}

    \item Wind power
    \begin{itemize}
        \item Advantage:
        \vspace{3cm}
        \item Disadvantage:
        \vspace{3cm}
    \end{itemize}

    \item Solar power
    \begin{itemize}
        \item Advantage:
        \vspace{3cm}
        \item Disadvantage:
        \vspace{3cm}
    \end{itemize}

    \item Fossil fuels (coal, oil or gas)
    \begin{itemize}
        \item Advantage:
        \vspace{3cm}
        \item Disadvantage:
        \vspace{3cm}
    \end{itemize}
\end{enumerate}

\newpage

\section*{3. Meeting Electricity Demand}

\textbf{Key ideas}

\begin{itemize}[noitemsep]
    \item Electricity demand changes throughout the day (e.g.\ higher in the evening).
    \item Power stations that run all the time provide the \textbf{base load}.
    \item Other power stations can be switched on and off quickly to meet \textbf{peak demand}.
    \item Some renewable resources are not always available (e.g.\ wind, solar).
\end{itemize}

\subsection*{Task 8: Base load and peak demand}

\begin{enumerate}[label=\arabic*)]
    \item What is meant by \textbf{base load}?
    \vspace{2cm}

    \item Which type of power station is usually used to provide base load in the UK? Explain why.
    \vspace{2.5cm}

    \item What is meant by \textbf{peak demand}? Give one example of a time when peak demand occurs.
    \vspace{2.5cm}

    \item Suggest two types of power station that are good for meeting peak demand and explain why.
    \vspace{3cm}
\end{enumerate}
\newpage
\subsection*{Task 9: Energy mix}

Countries usually use a \emph{mix} of energy resources.

\begin{enumerate}[label=\arabic*)]
    \item Give two reasons why it is a good idea to use a mixture of different energy resources instead of just one.
    \vspace{2.5cm}

    \item A country decides to shut down all of its nuclear power stations immediately and replace them only with wind turbines. Suggest two problems this might cause.
    \vspace{3cm}
\end{enumerate}

\newpage

\section*{4. Environmental and Social Issues}

\textbf{Key ideas}

\begin{itemize}[noitemsep]
    \item Different energy resources have different impacts on the environment and on people.
    \item Things to consider: air pollution, greenhouse gases, visual impact, noise, habitat destruction, fuel transport, waste.
\end{itemize}

\subsection*{Task 10: Matching statements}

Match each statement to the most suitable energy resource. You may use some resources more than once.

\textbf{Resource options:} coal, natural gas, nuclear, wind, solar, hydroelectric, tidal, biomass

\begin{enumerate}[label=\arabic*)]
    \item Produces long-lasting radioactive waste that must be stored safely.
    \vspace{0.5cm}
    \item Can cause flooding of large areas of land and destruction of habitats when built.
    \vspace{0.5cm}
    \item Produces carbon dioxide but can be turned on and off quickly.
    \vspace{0.5cm}
    \item No fuel costs and no carbon dioxide during operation, but depends on the weather.
    \vspace{0.5cm}
    \item Uses plant material or animal waste as fuel and can be considered carbon neutral.
    \vspace{0.5cm}
    \item Makes a humming noise, and some people think they spoil the view.
    \vspace{0.5cm}
    \item Produces sulfur dioxide and carbon dioxide, leading to acid rain and global warming.
    \vspace{1cm}
\end{enumerate}

\subsection*{Task 11: Short written questions}

Answer in full sentences.

\begin{enumerate}[label=\arabic*)]
    \item Why do many governments want to reduce the use of coal and oil for electricity generation?
    \vspace{2.5cm}

    \item Explain how using more renewable energy resources could help to slow down climate change.
    \vspace{2.5cm}

    \item Give one reason why some people are against building more wind farms, even though wind power is renewable.
    \vspace{2.5cm}

    \item Nuclear power does not produce carbon dioxide during operation, but some people are still worried about it. Explain why.
    \vspace{2.5cm}
\end{enumerate}

\newpage

\section*{5. Mixed Exam-Style Questions}

\subsection*{Task 12: Multiple choice}

Circle the correct answer for each question.

\begin{enumerate}[label=\arabic*)]
    \item Which of these is a \textbf{renewable} energy resource? \\
    A) Coal \hspace{0.7cm}
    B) Natural gas \hspace{0.7cm}
    C) Nuclear fuel \hspace{0.7cm}
    D) Solar

    \vspace{0.5cm}

    \item Which equation is correct? \\
    A) $P = E \times t$ \\
    B) $E = P \times t$ \\
    C) $E = \frac{t}{P}$ \\
    D) $P = \frac{t}{E}$

    \vspace{0.5cm}

    \item A power station that has a very short start-up time is most useful for: \\
    A) Providing base load \\
    B) Meeting peak demand \\
    C) Storing nuclear waste \\
    D) Reducing visual pollution

    \vspace{0.5cm}

    \item An appliance has an efficiency of 0.8. This means that: \\
    A) \SI{80}{\percent} of the input energy is wasted \\
    B) \SI{20}{\percent} of the input energy is useful \\
    C) \SI{80}{\percent} of the input energy is useful \\
    D) It does not waste any energy
\end{enumerate}

\vspace{0.5cm}
\newpage
\subsection*{Task 13: Data and calculations}

A small wind turbine has a power output of \SI{400}{\watt} when the wind is blowing at a steady speed.

\begin{enumerate}[label=\arabic*)]
    \item Calculate the energy transferred by the turbine in \SI{10}{\minute}. \\
    Show your working.
    \vspace{3cm}

    \item On a very windy day, the turbine produces \SI{2.0e6}{\joule} of energy in \SI{1}{\hour}. \\
    Calculate the average power output on that day.
    \vspace{3cm}

    \item The turbine transfers \SI{1.5e6}{\joule} of useful electrical energy in \SI{1}{\hour} when the total input energy from the wind is \SI{2.0e6}{\joule}. \\
    Calculate the efficiency as a percentage.
    \vspace{3cm}
\end{enumerate}

\newpage

\subsection*{Task 14: Longer written question}

Answer in full sentences and use key terms such as: \emph{renewable, non-renewable, base load, peak demand, reliability, carbon dioxide, pollution, cost}.

\textbf{Question:}

\emph{``A country wants to reduce its carbon dioxide emissions but still provide a reliable supply of electricity. Explain how using a mix of different energy resources (renewable and non-renewable) can help achieve this.''}

Use the space below to write a well-structured answer (aim for at least 8--10 lines).

\vspace{9cm}

\hrule
\vspace{0.3cm}

\textbf{Checklist (tick when done):}

\begin{itemize}[noitemsep]
    \item[$\square$] I have completed all calculation questions and shown my working.
    \item[$\square$] I have explained the key ideas of renewable and non-renewable resources.
    \item[$\square$] I have compared advantages and disadvantages of different energy resources.
    \item[$\square$] I have answered the longer written question using key scientific terms.
\end{itemize}

\vfill

\end{document}

